\documentclass[../../../include/open-logic-section]{subfiles}

\begin{document}

\olfileid{sth}{ord-arithmetic}{add}
\olsection{الجمع الترتيبي}

لنفرض أننا نريد جمع~$\OLId{greek-variable}{alpha}$ و$\beta$. نضع نسخة من~$\beta$ عقب
نسخة من~$\OLId{greek-variable}{alpha}$ مباشرة. وإنما نأخذ \emph{نسختين} لأن
\olref[ordinals][basic]{ordinalsaresubsets} يبين أن إما
$\OLId{greek-variable}{alpha} \subseteq \beta$ أو~$\beta \subseteq \OLId{greek-variable}{alpha}$. ومعنى هذا في
الحدس أن نجتاز خطوات نمطها~$\OLId{greek-variable}{alpha}$، ثم نجتاز خطوات نمطها~$\beta$.
والترتيب الناتج جيد، فيتماثل، بموجب
\olref[ordinals][ordtype]{thmOrdinalRepresentation}، مع عدد ترتيبي
واحد. وذلك العدد هو الذي نسميه \emph{مجموع}~$\OLId{greek-variable}{alpha}$ و$\beta$.

هذه جملة الفكرة وراء الجمع الترتيبي. ولضبطها نبدأ بمعنى أخذ
\emph{نسخ} من المجموعات: نستعمل وسمين اعتباطيين، هما~$\OLMathNumeral{0}$ و$\OLMathNumeral{1}$، لنعلم
من أين جاء كل شيء.

\begin{defn}\ollabel{defdissum}
\emph{الجمع المنفصل} لـ$\OLId{latin-looped}{A}$ و$\OLId{latin-looped}{B}$ هو
$\OLId{latin-looped}{A} \disjointsum \OLId{latin-looped}{B} = (\OLId{latin-looped}{A}\times \{\OLMathNumeral{0}\}) \cup (\OLId{latin-looped}{B} \times \{\OLMathNumeral{1}\})$.
\end{defn}

ثم نضع ترتيبًا على أزواج الأعداد الترتيبية.

\begin{defn}
لأي أعداد ترتيبية~$\OLId{greek-variable}{alpha}_\OLMathNumeral{1}, \OLId{greek-variable}{alpha}_\OLMathNumeral{2}, \beta_\OLMathNumeral{1}, \beta_\OLMathNumeral{2}$، نقول إن:
\begin{align*}
	\tuple{\OLId{greek-variable}{alpha}_\OLMathNumeral{1}, \OLId{greek-variable}{alpha}_\OLMathNumeral{2}} \rlexless \tuple{\beta_\OLMathNumeral{1}, \beta_\OLMathNumeral{2}}\text{ إذا وفقط إذا كان }&
	\text{إما $\alpha_2 \in \beta_2$،}\\
	& \text{أو كان معًا $\alpha_2 = \beta_2$ و$\alpha_1 \in \beta_1$}
\end{align*}
\end{defn}

وهذا ترتيب \emph{معجمي معكوس}؛ إذ يبدأ بالعنصر الثاني، ثم يرجع إلى
الأول. وقد أردنا أن يكون~$\OLId{greek-variable}{alpha} \ordplus \beta$ نمط ترتيب نسخة من
$\OLId{greek-variable}{alpha}$ تأتي بعدها نسخة من~$\beta$، فنقول:

\begin{defn}\ollabel{defordplus}
لأي عددين ترتيبيين~$\OLId{greek-variable}{alpha}$ و$\beta$، يكون مجموعهما
$\OLId{greek-variable}{alpha} \ordplus \beta =
\ordtype{\OLId{greek-variable}{alpha} \disjointsum \beta, \rlexless}$.
\end{defn}
\noindent
وفي هذا تسامح يسير في الرمز؛ فحقنا أن نكتب
«$\Setabs{\tuple{\OLId{latin-ordinary}{x},\OLId{latin-ordinary}{y}}\in(\OLId{greek-variable}{alpha}\disjointsum\beta)^\OLMathNumeral{2}}{\OLId{latin-ordinary}{x} \rlexless \OLId{latin-ordinary}{y}}$»
مكان «$\rlexless$». وسنبقي هذا التسامح طلبًا للاختصار.

وتدل النتيجة الآتية، مع
\olref[ordinals][ordtype]{thmOrdinalRepresentation}، على سلامة
التعريف:

\begin{lem}\ollabel{ordsumlessiswo}
$\tuple{\OLId{greek-variable}{alpha} \disjointsum \beta, \rlexless}$ ترتيب جيد، لأي عددين
ترتيبيين~$\OLId{greek-variable}{alpha}$ و$\beta$.
\end{lem}

\begin{proof}
اتصال~$\rlexless$ على~$\OLId{greek-variable}{alpha} \disjointsum \beta$ بيّن. ولإثبات حسن
التأسيس، ثبت مجموعة غير خالية~$\OLId{latin-looped}{X} \subseteq \OLId{greek-variable}{alpha} \disjointsum \beta$.
ولتكن~$\OLId{latin-looped}{Y}$ مجموعة عناصر~$\OLId{latin-looped}{X}$ التي إحداثيها الثاني أصغر ما يمكن، أي
$\OLId{latin-looped}{Y} = \Setabs{\tuple{\gamma, \OLId{latin-ordinary}{i}} \in \OLId{latin-looped}{X}}{(\forall \tuple{\delta, \OLId{latin-ordinary}{j}} \in \OLId{latin-looped}{X})\OLId{latin-ordinary}{i} \leq \OLId{latin-ordinary}{j}}$.
ثم خذ من~$\OLId{latin-looped}{Y}$ العنصر الذي إحداثيه الأول أصغر.
\end{proof}
\noindent
فقد حصلنا على تعريف صريح حسن للجمع الترتيبي. ومن خواصه القريبة، مع
تذكر~$\OLMathNumeral{1} = \{\OLMathNumeral{0}\}$ من~\olref[z][infinity-again]{defnomega}، ما يأتي:
\begin{prop}
$\OLId{greek-variable}{alpha} \ordplus \OLMathNumeral{1} = \ordsucc{\OLId{greek-variable}{alpha}}$، لأي عدد ترتيبي~$\OLId{greek-variable}{alpha}$.
\end{prop}

\begin{proof}
اعتبر التماثل~$\OLId{latin-ordinary}{f}$ من
$\ordsucc{\OLId{greek-variable}{alpha}} = \OLId{greek-variable}{alpha} \cup \{\OLId{greek-variable}{alpha}\}$ إلى
$\OLId{greek-variable}{alpha}\disjointsum\OLMathNumeral{1} = (\OLId{greek-variable}{alpha} \times \{\OLMathNumeral{0}\})
\cup (\{\OLMathNumeral{0}\} \times \{\OLMathNumeral{1}\})$، المعطى بـ
$\OLId{latin-ordinary}{f}(\gamma) = \tuple{\gamma, \OLMathNumeral{0}}$ لكل~$\gamma \in \OLId{greek-variable}{alpha}$، وبـ
$\OLId{latin-ordinary}{f}(\OLId{greek-variable}{alpha}) = \tuple{\OLMathNumeral{0}, \OLMathNumeral{1}}$.
\end{proof}
\noindent
وللجمع كذلك أوصاف عودية يسهل إثباتها.

\begin{lem}\ollabel{ordadditionrecursion}
لأي عددين ترتيبيين~$\OLId{greek-variable}{alpha}, \beta$، لدينا:
\begin{align*}
	\OLId{greek-variable}{alpha}\ordplus \OLMathNumeral{0} &= \OLId{greek-variable}{alpha}\\
	\OLId{greek-variable}{alpha} \ordplus  (\beta\ordplus \OLMathNumeral{1}) &= (\OLId{greek-variable}{alpha} \ordplus  \beta) \ordplus  \OLMathNumeral{1}\\
	\OLId{greek-variable}{alpha}  \ordplus  \beta &= \supstrict_{\delta < \beta}(\OLId{greek-variable}{alpha} \ordplus  \delta) && \text{إذا كان $\beta$ عددًا ترتيبيًا حديًا}
\end{align*}
\end{lem}

\begin{proof}
ننظر في الحالات على ترتيبها. أولًا:
\begin{align*}
	\OLId{greek-variable}{alpha} \ordplus  \OLMathNumeral{0}
	&= \ordtype{\OLId{greek-variable}{alpha} \times \{\OLMathNumeral{0}\}, \rlexless}\\
	&= \OLId{greek-variable}{alpha}\\
	\OLId{greek-variable}{alpha} \ordplus (\beta \ordplus  \OLMathNumeral{1})
	%&= \ordtype{(\alpha\times \{0\}) \cup (\ordtype{\beta \ordplus  1}\times \{1\}), \rlexless} \\
	&= \ordtype{(\OLId{greek-variable}{alpha}\times \{\OLMathNumeral{0}\}) \cup (\ordsucc{\beta}\times \{\OLMathNumeral{1}\}), \rlexless} \\
	&= \ordtype{(\OLId{greek-variable}{alpha}\times \{\OLMathNumeral{0}\}) \cup (\beta \times \{\OLMathNumeral{1}\}), \rlexless} \ordplus \OLMathNumeral{1}\\
	&= (\OLId{greek-variable}{alpha} \ordplus  \beta) \ordplus  \OLMathNumeral{1}
\end{align*}
والآن لتكن~$\beta \neq \emptyset$ حدية. متى كان~$\delta < \beta$،
كان~$\delta\ordplus \OLMathNumeral{1} < \beta$ أيضًا، ومن ثم كان
$\OLId{greek-variable}{alpha} \ordplus \delta$ مقطعًا ابتدائيًا حقيقيًا من
$\OLId{greek-variable}{alpha} \ordplus \beta$. فيكون~$\OLId{greek-variable}{alpha} \ordplus \beta$ حدًا أعلى
صارمًا للمجموعة
$\OLId{latin-looped}{X} = \Setabs{\OLId{greek-variable}{alpha} \ordplus \delta}{\delta < \beta}$. ثم إذا كان
$\OLId{greek-variable}{alpha} \leq \gamma < \OLId{greek-variable}{alpha} \ordplus \beta$، كان بالضرورة
$\gamma = \OLId{greek-variable}{alpha} \ordplus \delta$ لبعض~$\delta < \beta$. ولهذا
$\OLId{greek-variable}{alpha} \ordplus \beta =
\supstrict_{\delta< \beta}(\OLId{greek-variable}{alpha}\ordplus \delta)$.
\end{proof}

غير أن في الأمر وجهًا آخر. فقد كان في وسعنا أن نعرّف الجمع الترتيبي
بـ\emph{مجرد} مبرهنة العودية المتناهية التجاوز، فنقرر معادلات
العودية الواردة في~\olref{ordadditionrecursion}، مع استعمال
«$\ordsucc{\beta}$» مكان «$\beta \ordplus \OLMathNumeral{1}$».

فللعمليات على الأعداد الترتيبية طريقان: طريق \emph{تركيبي} نبني فيه
مجموعة مرتبة ترتيبًا جيدًا ثم نأخذ نمط ترتيبها، وطريق \emph{عودي}
نقرر فيه معادلات العودية. وإذا أحكم العمل اتفق الطريقان؛ لأن مبرهنة
العودية المتناهية التجاوز تضمن أن تلك المعادلات تعين أعدادًا ترتيبية
\emph{فريدة}.

والحساب الترتيبي يشبه، من وجوه كثيرة، حساب الأعداد الطبيعية؛ ومن ذلك:

\begin{lem}\ollabel{ordinaladditionisnice}
إذا كانت~$\OLId{greek-variable}{alpha}, \beta, \gamma$ أعدادًا ترتيبية، فإن:
\begin{enumerate}
	\item\ollabel{ordaddition1} إذا كان~$\beta < \gamma$، فإن
	$\OLId{greek-variable}{alpha} \ordplus \beta < \OLId{greek-variable}{alpha} \ordplus \gamma$؛
	\item\ollabel{ordaddition2} إذا كان~$\OLId{greek-variable}{alpha} \ordplus \beta =
	\OLId{greek-variable}{alpha}\ordplus \gamma$، فإن~$\beta = \gamma$؛
	\item\ollabel{ordaddition3} $\OLId{greek-variable}{alpha} \ordplus (\beta \ordplus
	\gamma) = (\OLId{greek-variable}{alpha} \ordplus \beta) \ordplus \gamma$، أي إن الجمع
	تجميعي؛
	\item\ollabel{ordaddition4} إذا كان~$\OLId{greek-variable}{alpha} \leq \beta$، فإن
	$\OLId{greek-variable}{alpha} \ordplus \gamma \leq \beta \ordplus \gamma$.
\end{enumerate}
\end{lem}

\begin{proof}
نبرهن~\olref{ordaddition3} ونكل الباقي إلى التمرين. نستعمل الاستقراء
البسيط المتناهي التجاوز على~$\gamma$، مع
\olref{ordadditionrecursion}. فإذا~$\gamma = \OLMathNumeral{0}$:
\[
(\OLId{greek-variable}{alpha} \ordplus  \beta) \ordplus  \OLMathNumeral{0} = \OLId{greek-variable}{alpha} \ordplus  \beta  = \OLId{greek-variable}{alpha} \ordplus  (\beta \ordplus  \OLMathNumeral{0})
\]
وإذا~$\gamma = \delta\ordplus \OLMathNumeral{1}$، فافترض استقرائيًا أن
$(\OLId{greek-variable}{alpha} \ordplus  \beta) \ordplus  \delta =
\OLId{greek-variable}{alpha} \ordplus  (\beta \ordplus \delta)$؛ ثم استعمل
\olref{ordadditionrecursion} ثلاث مرات:
\begin{align*}
	(\OLId{greek-variable}{alpha} \ordplus  \beta) \ordplus  (\delta \ordplus  \OLMathNumeral{1}) & = ((\OLId{greek-variable}{alpha} \ordplus  \beta) \ordplus  \delta)\ordplus \OLMathNumeral{1}\\
	& = (\OLId{greek-variable}{alpha} \ordplus  (\beta \ordplus  \delta)) \ordplus  \OLMathNumeral{1}\\
	& = \OLId{greek-variable}{alpha} \ordplus  ((\beta \ordplus  \delta)\ordplus \OLMathNumeral{1})\\
	& = \OLId{greek-variable}{alpha} \ordplus  (\beta \ordplus  (\delta\ordplus \OLMathNumeral{1}))
\end{align*}	
وأما إذا كانت~$\gamma$ عددًا ترتيبيًا حديًا، فافترض استقرائيًا أنه
متى كان~$\delta \in \gamma$، فإن
$(\OLId{greek-variable}{alpha} \ordplus  \beta) \ordplus  \delta =
\OLId{greek-variable}{alpha} \ordplus  (\beta \ordplus  \delta)$؛ وعندئذ:
\begin{align*}
	(\OLId{greek-variable}{alpha} \ordplus  \beta) \ordplus  \gamma & = \supstrict_{\delta < \gamma}((\OLId{greek-variable}{alpha} \ordplus  \beta) \ordplus  \delta) \\
	&= \supstrict_{\delta < \gamma}(\OLId{greek-variable}{alpha} \ordplus  (\beta \ordplus  \delta))\\
	&= \OLId{greek-variable}{alpha} \ordplus  \supstrict_{\delta < \gamma}(\beta \ordplus  \delta)\\
	& = \OLId{greek-variable}{alpha} \ordplus  (\beta \ordplus  \gamma)
\end{align*}
\end{proof}

\begin{prob}
أثبت بقية~\olref[sth][ord-arithmetic][add]{ordinaladditionisnice}.
\end{prob}

بهذه الخواص يبدو الجمع الترتيبي مألوفًا، غير أنه يفارق جمع الأعداد
الطبيعية في وجه جوهري.

\begin{prop}\ollabel{ordsumnotcommute}
الجمع الترتيبي {ليس} تبادليًا؛
$\OLMathNumeral{1} \ordplus \omega = \omega < \omega \ordplus \OLMathNumeral{1}$.
\end{prop}

\begin{proof}
لاحظ أن~$\OLMathNumeral{1} \ordplus \omega = \supstrict_{\OLId{latin-ordinary}{n} < \omega} (\OLMathNumeral{1} \ordplus \OLId{latin-ordinary}{n})
= \omega \in \omega \cup \{\omega\} = \ordsucc{\omega} = \omega
\ordplus \OLMathNumeral{1}$.
\end{proof}
\noindent
وقد يفجأ هذا أول النظر، ولكنه لا ينبغي أن يفجأ. فإن
$\OLMathNumeral{1} \ordplus \omega$ هو نمط الترتيب الحاصل بإدخال عنصر زائد
\emph{قبل} الأعداد الطبيعية كلها. وكما في حجة فندق هيلبرت في
\olref[sfr][infinite][hilbert]{sec}، لا يغير ذلك العنصر الأول نمط
الترتيب الكلي في الحدس. وأما~$\omega \ordplus \OLMathNumeral{1}$ فهو نمط الترتيب
الحاصل بإلحاق عنصر \emph{بعد} الأعداد الطبيعية كلها؛ وبين الحالين
فرق بين.

\end{document}
